% Major pandemics 1918 -> 2020 with mortality bars.
% Vol VI Ch 8 §8.7.
\begin{figure}[htbp]
\centering
\begin{tikzpicture}
\begin{axis}[
  width=13cm,
  height=7cm,
  xlabel={Year},
  ylabel={Estimated global mortality},
  xmin=1910, xmax=2030,
  ymin=0, ymax=60000000,
  xtick={1918, 1957, 1968, 1977, 2003, 2009, 2020},
  xticklabel style={rotate=45, anchor=north east, font=\small},
  ytick={0, 1000000, 7000000, 50000000},
  yticklabels={$0$, $1\,\text{M}$, $7\,\text{M}$, $50\,\text{M}$},
  ymode=log,
  ymin=100, ymax=1e8,
  ybar,
  bar width=10pt,
  grid=major,
  grid style={draw=gray, very thin, opacity=0.3},
]
  \addplot[fill=engamber!60, draw=engamber] coordinates {
    (1918, 50000000)
    (1957, 1100000)
    (1968, 1000000)
    (1977, 700000)
    (2003, 774)
    (2009, 360000)
    (2020, 7000000)
  };
\end{axis}
\end{tikzpicture}
\caption[Pandemic mortality timeline]{Estimated global mortality
of major influenza pandemics and respiratory-pathogen outbreaks
since 1918 (log scale). The 1918--1920 H1N1 pandemic killed an
estimated 50 million globally
\cite{paper:v6c08-johnson-mueller-1918-flu-2002}. Subsequent
influenza pandemics (1957 H2N2, 1968 H3N2, 1977 H1N1) and
respiratory-pathogen outbreaks (2003 SARS-CoV-1, 2009 H1N1, 2020
SARS-CoV-2) killed 770--7 million each. The figures for 2009
match the Dawood et al.\ 2012 estimate
\cite{paper:v6c08-dawood-h1n1-2012}; 2020 follows the WHO
excess-mortality estimate.}
\label{fig:vol06:ch08:pandemic-timeline}
\end{figure}
